\chapter{Vetores}\label{ch:g10:vectors}

Um \omterm{def:g10:vectors:vector}{vetor} codifica um deslocamento: uma direção com um sentido e um
comprimento, independentemente do ponto de partida. Os \omterm{def:g10:vectors:vector}{vetores} dão
demonstrações limpas de fatos geométricos e, uma vez que as \omterm{def:g10:coordgeom:system}{coordenadas}
entram em cena, reduzem-nos a pequenos cálculos. Serão uma ferramenta
central no \cref{ch:g10:lines} e, mais tarde, para o produto escalar no
\cref{ch:g11:scal}.

\section{Translações e vetores}

\begin{definition}[Vetor]\label{def:g10:vectors:vector}
Dados dois pontos $A$ e $B$, o \emph{vetor}\index{vetor} $\vect{AB}$
representa o deslocamento de $A$ até $B$: ele tem uma \emph{direção com
sentido} (a direção da reta $(AB)$, junto com o modo de percorrê-la, de
$A$ rumo a $B$) e um \emph{comprimento} (a distância $AB$, também escrita
$\norm{\vect{AB}}$). Dois vetores são \emph{iguais} quando codificam o
mesmo deslocamento: $\vect{AB} = \vect{CD}$ significa que aplicar os dois
deslocamentos a um ponto qualquer dá o mesmo resultado.
\end{definition}

\begin{proposition}[Vetores iguais e paralelogramos]\label{prop:g10:vectors:parallelogram}
$\vect{AB} = \vect{CD}$ exatamente quando $ABDC$ (nesta ordem!) é um
paralelogramo, possivelmente achatado --- ou, o que é equivalente, quando
$[AD]$ e $[BC]$ têm o mesmo \omterm{prop:g10:coordgeom:midpoint}{ponto médio}.
\end{proposition}
\admitted

\begin{omfigure}
\begin{tikzpicture}[scale=1.0]
  \coordinate (A) at (0,0);
  \coordinate (B) at (2.2,0.8);
  \coordinate (C) at (1.0,-1.2);
  \coordinate (D) at (3.2,-0.4);
  \draw[omExo, thin] (A) -- (C);
  \draw[omExo, thin] (B) -- (D);
  \draw[-Stealth, omDef, thick] (A) -- (B)
    node[midway, above, font=\small] {$\vect{AB}$};
  \draw[-Stealth, omDef, thick] (C) -- (D)
    node[midway, below, font=\small] {$\vect{CD}$};
  \fill (A) circle (1.5pt) node[left, font=\small] {$A$};
  \fill (B) circle (1.5pt) node[right, font=\small] {$B$};
  \fill (C) circle (1.5pt) node[left, font=\small] {$C$};
  \fill (D) circle (1.5pt) node[right, font=\small] {$D$};
\end{tikzpicture}

{\small \omterm{def:g10:vectors:vector}{Vetores} iguais $\vect{AB} = \vect{CD}$: mesma direção, mesmo
sentido, mesmo comprimento. O quadrilátero $ABDC$ é um paralelogramo.}
\end{omfigure}

\begin{notation}
Um \omterm{def:g10:vectors:vector}{vetor} costuma ser designado por uma única letra, $\vec u$ ou $\vec v$,
quando as extremidades não importam. O \emph{\omterm{def:g10:vectors:vector}{vetor} nulo}
$\vec 0 = \vect{AA}$ é o deslocamento que não move nada. O
\emph{oposto} de $\vect{AB}$ é $\vect{BA}$: mesmo comprimento, sentido
contrário. Escrevemos $\vect{BA} = -\vect{AB}$.
\end{notation}

\section{Somar vetores}

\begin{definition}[Soma de dois vetores]\label{def:g10:vectors:sum}
A \emph{soma}\index{vetor!soma} $\vec u + \vec v$ é o deslocamento obtido
efetuando $\vec u$ e em seguida $\vec v$.
\end{definition}

\begin{theorem}[Relação de Chasles]\label{thm:g10:vectors:chasles}
Para quaisquer três pontos $A$, $B$, $C$:
\[
\vect{AB} + \vect{BC} = \vect{AC}.
\]
\end{theorem}

\begin{proof}
Ir de $A$ a $B$ e depois de $B$ a $C$ é um deslocamento que leva $A$ até
$C$: é o deslocamento codificado por $\vect{AC}$.
\end{proof}

\begin{remark}[Regra do paralelogramo]
Para somar dois \omterm{def:g10:vectors:vector}{vetores} traçados a partir do mesmo ponto,
$\vect{AB} + \vect{AC}$, complete o paralelogramo $ABDC$: a \omterm{def:g10:vectors:sum}{soma} é a
diagonal $\vect{AD}$. De fato, $\vect{AC} = \vect{BD}$, logo
$\vect{AB} + \vect{AC} = \vect{AB} + \vect{BD} = \vect{AD}$ pela relação
de Chasles.
\end{remark}

\begin{omfigure}
\begin{tikzpicture}[scale=1.05]
  \coordinate (A) at (0,0);
  \coordinate (B) at (2.4,0.5);
  \coordinate (C) at (3.4,2.0);
  \draw[-Stealth, omDef, thick] (A) -- (B)
    node[midway, below, font=\small] {$\vec u$};
  \draw[-Stealth, omThm, thick] (B) -- (C)
    node[midway, right, font=\small] {$\vec v$};
  \draw[-Stealth, omMeth, thick] (A) -- (C)
    node[midway, above left, font=\small] {$\vec u + \vec v$};
\end{tikzpicture}\hspace{1.4cm}%
\begin{tikzpicture}[scale=1.05]
  \coordinate (A) at (0,0);
  \coordinate (B) at (2.4,0.5);
  \coordinate (C) at (1.0,1.5);
  \coordinate (D) at (3.4,2.0);
  \draw[omExo, thin] (B) -- (D) -- (C);
  \draw[-Stealth, omDef, thick] (A) -- (B)
    node[midway, below, font=\small] {$\vect{AB}$};
  \draw[-Stealth, omThm, thick] (A) -- (C)
    node[midway, left, font=\small] {$\vect{AC}$};
  \draw[-Stealth, omMeth, thick] (A) -- (D)
    node[pos=0.75, below right, font=\small] {$\vect{AD}$};
  \fill (D) circle (1.2pt) node[right, font=\small] {$D$};
\end{tikzpicture}

{\small Duas visões da \omterm{def:g10:vectors:sum}{soma}: ponta com origem (relação de Chasles, à
esquerda) e a regra do paralelogramo (à direita).}
\end{omfigure}

\begin{example}[Simplificar com Chasles]\label{ex:g10:vectors:chasles}
Simplifique $\vect{MN} + \vect{NP} + \vect{PQ}$: encadear os
deslocamentos dá $\vect{MQ}$. Simplifique $\vect{AB} - \vect{AC}$:
reescreva a diferença como \omterm{def:g10:vectors:sum}{soma} com o \omterm{def:g10:vectors:vector}{vetor} oposto,
\[
\vect{AB} - \vect{AC} = \vect{AB} + \vect{CA}
= \vect{CA} + \vect{AB} = \vect{CB}.
\]
\end{example}

\section{Multiplicar um vetor por um número}

\begin{definition}[Múltiplo escalar]\label{def:g10:vectors:scalar}
Seja $\vec u$ um \omterm{def:g10:vectors:vector}{vetor} não nulo e $k$ um \omterm{def:g10:numbers:sets}{número real}. O \omterm{def:g10:vectors:vector}{vetor}
$k\vec u$\index{vetor!múltiplo escalar} tem a mesma direção de $\vec u$,
comprimento $\abs{k} \times \norm{\vec u}$ e aponta no mesmo sentido de
$\vec u$ se $k > 0$, no sentido contrário se $k < 0$. Para $k = 0$,
$0\vec u = \vec 0$.
\end{definition}

\begin{definition}[Colinearidade]\label{def:g10:vectors:collinear}
Dois \omterm{def:g10:vectors:vector}{vetores} $\vec u$ e $\vec v$ são
\emph{colineares}\index{vetores colineares} quando têm a mesma direção,
isto é, quando $\vec v = k \vec u$ para algum real $k$ (ou quando um
deles é $\vec 0$).
\end{definition}

\begin{remark}
A colinearidade é a linguagem vetorial do paralelismo e do alinhamento:
\begin{itemize}
  \item as retas $(AB)$ e $(CD)$ são paralelas exatamente quando
        $\vect{AB}$ e $\vect{CD}$ são \omterm{def:g10:vectors:collinear}{colineares};
  \item os pontos $A$, $B$, $C$ são \omterm{def:g10:vectors:collinear}{colineares} exatamente quando
        $\vect{AB}$ e $\vect{AC}$ o são.
\end{itemize}
\end{remark}

\section{Coordenadas de vetores}

\begin{definition}[Coordenadas de um vetor]\label{def:g10:vectors:coords}
Em um \omterm{def:g10:coordgeom:system}{sistema de coordenadas}, as \emph{\omterm{def:g10:coordgeom:system}{coordenadas}} do \omterm{def:g10:vectors:vector}{vetor} $\vect{AB}$
são os números que descrevem o deslocamento:
\[
\vect{AB}\,(x_B - x_A,\ y_B - y_A).
\]
Um \omterm{def:g10:vectors:vector}{vetor} $\vec u\,(a, b)$ move todo ponto $a$ unidades na horizontal e
$b$ unidades na vertical.
\end{definition}

\begin{proposition}[Calcular com coordenadas]\label{prop:g10:vectors:coordrules}
Sejam $\vec u\,(a, b)$ e $\vec v\,(c, d)$, e seja $k \in \R$.
\begin{enumerate}
  \item $\vec u = \vec v$ exatamente quando $a = c$ e $b = d$;
  \item $\vec u + \vec v$ tem \omterm{def:g10:coordgeom:system}{coordenadas} $(a + c,\ b + d)$;
  \item $k\vec u$ tem \omterm{def:g10:coordgeom:system}{coordenadas} $(ka,\ kb)$;
  \item $\norm{\vec u} = \sqrt{a^2 + b^2}$ (\omterm{def:g10:coordgeom:system}{sistema ortonormal}).
\end{enumerate}
\end{proposition}

\begin{proof}
1--3 reenunciam, coordenada a coordenada, o que os deslocamentos fazem:
por exemplo, efetuar $\vec u$ e depois $\vec v$ move um ponto
horizontalmente de $a$ e depois de $c$, portanto de $a + c$ ao todo. O
item 4 é a fórmula da distância do \cref{thm:g10:coordgeom:distance}
aplicada a um segmento que representa $\vec u$.
\end{proof}

\begin{theorem}[Critério de colinearidade]\label{thm:g10:vectors:det}
Dois \omterm{def:g10:vectors:vector}{vetores} $\vec u\,(a, b)$ e $\vec v\,(c, d)$ são \omterm{def:g10:vectors:collinear}{colineares} se, e
somente se,
\[
ad - bc = 0 .
\]
\end{theorem}

\begin{proof}
Se $\vec v = k\vec u$, então $c = ka$ e $d = kb$, logo
$ad - bc = a(kb) - b(ka) = 0$. O mesmo vale se $\vec u = k \vec v$ ou se
um dos \omterm{def:g10:vectors:vector}{vetores} é nulo.

Reciprocamente, suponha $ad - bc = 0$ com $\vec u \neq \vec 0$, digamos
$a \neq 0$ (o caso $b \neq 0$ é análogo). Ponha $k = \frac{c}{a}$; então
$c = ka$, e $ad = bc = b(ka)$ dá, após dividir por $a \neq 0$, $d = kb$.
Logo $\vec v = k\vec u$.
\end{proof}

\begin{example}\label{ex:g10:vectors:det}
Os \omterm{def:g10:vectors:vector}{vetores} $\vec u\,(3, -2)$ e $\vec v\,(-6, 4)$ são \omterm{def:g10:vectors:collinear}{colineares}? Calcule
$ad - bc = 3 \times 4 - (-2)\times(-6) = 12 - 12 = 0$: sim, e de fato
$\vec v = -2\vec u$. Já para $\vec u\,(3, -2)$ e $\vec w\,(1, 5)$:
$3 \times 5 - (-2) \times 1 = 17 \neq 0$: não são \omterm{def:g10:vectors:collinear}{colineares}.
\end{example}

\begin{method}[Colinearidade de pontos e paralelismo]\label{met:g10:vectors:alignment}
Para provar que três pontos $A$, $B$, $C$ são \omterm{def:g10:vectors:collinear}{colineares}:
\begin{enumerate}
  \item calcule as \omterm{def:g10:coordgeom:system}{coordenadas} de $\vect{AB}$ e de $\vect{AC}$;
  \item verifique o critério $ad - bc = 0$ do
        \cref{thm:g10:vectors:det};
  \item conclua: os \omterm{def:g10:vectors:vector}{vetores} são \omterm{def:g10:vectors:collinear}{colineares} e partilham o ponto $A$, logo
        os três pontos são \omterm{def:g10:vectors:collinear}{colineares}.
\end{enumerate}
O mesmo cálculo com $\vect{AB}$ e $\vect{CD}$ prova que as retas $(AB)$ e
$(CD)$ são paralelas.
\end{method}

\begin{example}\label{ex:g10:vectors:alignment}
Sejam $A(1, 2)$, $B(3, 5)$ e $C(7, 11)$. Então $\vect{AB}\,(2, 3)$ e
$\vect{AC}\,(6, 9)$, e $2 \times 9 - 3 \times 6 = 18 - 18 = 0$: os pontos
são \omterm{def:g10:vectors:collinear}{colineares} (na verdade, $\vect{AC} = 3\vect{AB}$).
\end{example}

\begin{omfigure}
\begin{tikzpicture}[scale=0.62]
  \draw[thin, gray!40] (-0.4,-0.4) grid (8.4,5.4);
  \draw[-Stealth] (-0.6,0) -- (8.6,0);
  \draw[-Stealth] (0,-0.6) -- (0,5.6);
  \coordinate (A) at (1,1);
  \coordinate (B) at (3,2);
  \coordinate (C) at (7,4);
  \draw[-Stealth, omDef, very thick] (A) -- (B)
    node[midway, above left=-2pt, font=\small] {$\vect{AB}$};
  \draw[-Stealth, omThm, thick] (A) -- (C)
    node[pos=0.8, below right=-2pt, font=\small] {$\vect{AC} = 3\vect{AB}$};
  \fill (A) circle (2pt) node[above left, font=\small] {$A$};
  \fill (B) circle (2pt) node[above left, font=\small] {$B$};
  \fill (C) circle (2pt) node[above left, font=\small] {$C$};
\end{tikzpicture}

{\small Pontos \omterm{def:g10:vectors:collinear}{colineares}: $\vect{AC}$ é um múltiplo escalar de
$\vect{AB}$.}
\end{omfigure}

\section{Exercícios}

\begin{exercise}[$\star$]\label{exo:g10:vectors:1}
Simplifique usando a relação de Chasles:
\[
\vect{AB} + \vect{BD}, \qquad
\vect{KL} + \vect{LM} + \vect{MK}, \qquad
\vect{AC} - \vect{BC}, \qquad
\vect{AB} + \vect{CA}.
\]
\end{exercise}

\begin{exercise}[$\star$]\label{exo:g10:vectors:2}
Sejam $A(2, 1)$, $B(5, 3)$, $C(-1, 4)$. Calcule as \omterm{def:g10:coordgeom:system}{coordenadas} de
$\vect{AB}$, $\vect{BC}$, $\vect{AC}$ e verifique nas \omterm{def:g10:coordgeom:system}{coordenadas} que
$\vect{AB} + \vect{BC} = \vect{AC}$.
\end{exercise}

\begin{exercise}[$\star$]\label{exo:g10:vectors:3}
Sejam $\vec u\,(2, -3)$ e $\vec v\,(-1, 4)$. Dê as \omterm{def:g10:coordgeom:system}{coordenadas} de
$\vec u + \vec v$, $3\vec u$, $2\vec u - \vec v$ e calcule
$\norm{\vec u}$.
\end{exercise}

\begin{exercise}[$\star$]\label{exo:g10:vectors:4}
Em cada caso, diga se $\vec u$ e $\vec v$ são \omterm{def:g10:vectors:collinear}{colineares}:
\[
\text{(a) } \vec u\,(4, 6),\ \vec v\,(6, 9);
\qquad
\text{(b) } \vec u\,(2, -5),\ \vec v\,(-4, 10);
\qquad
\text{(c) } \vec u\,(3, 1),\ \vec v\,(1, 3).
\]
\end{exercise}

\begin{exercise}[$\star$]\label{exo:g10:vectors:5}
Sejam $A(0, 2)$, $B(4, 0)$, $C(6, 3)$. Encontre as \omterm{def:g10:coordgeom:system}{coordenadas} do ponto
$D$ tal que $ABCD$ seja um paralelogramo, isto é, tal que
$\vect{AB} = \vect{DC}$.
\end{exercise}

\begin{exercise}[$\star\star$]\label{exo:g10:vectors:6}
Sejam $A(-2, 1)$, $B(1, 3)$, $C(4, -1)$.
\begin{enumerate}
  \item Calcule as \omterm{def:g10:coordgeom:system}{coordenadas} do ponto $M$ tal que
        $\vect{AM} = \vect{AB} + \vect{AC}$.
  \item Que quadrilátero é $ABMC$? Justifique.
\end{enumerate}
\end{exercise}

\begin{exercise}[$\star\star$]\label{exo:g10:vectors:7}
Sejam $A(1, -1)$, $B(4, 1)$, $C(10, 5)$. Os pontos $A$, $B$, $C$ são
\omterm{def:g10:vectors:collinear}{colineares}? Mesma pergunta para $A(0, 3)$, $B(2, 2)$, $C(8, -1)$.
\end{exercise}

\begin{exercise}[$\star\star$]\label{exo:g10:vectors:8}
Sejam $A(1, 2)$, $B(5, 4)$, $C(6, 1)$, $D(2, -1)$. Mostre que $(AB)$ e
$(CD)$ são paralelas e, em seguida, que $ABCD$ é um paralelogramo. Qual
das duas verificações implica a outra?
\end{exercise}

\begin{exercise}[$\star\star$]\label{exo:g10:vectors:9}
Seja $ABC$ um triângulo e seja $I$ o ponto definido por
$\vect{AI} = \frac23 \vect{AB}$. Com $A(0,0)$, $B(6, 3)$, $C(2, 5)$
(para que as contas fiquem simples):
\begin{enumerate}
  \item calcule as \omterm{def:g10:coordgeom:system}{coordenadas} de $I$;
  \item seja $J$ tal que $\vect{CJ} = \frac23\vect{CB}$; calcule as
        \omterm{def:g10:coordgeom:system}{coordenadas} de $J$;
  \item mostre que $(IJ)$ é paralela a $(AC)$.
\end{enumerate}
\end{exercise}

\begin{exercise}[$\star\star\star$]\label{exo:g10:vectors:10}
Seja $ABCD$ um quadrilátero qualquer e sejam $I$, $J$, $K$, $L$ os
\omterm{prop:g10:coordgeom:midpoint}{pontos médios} de $[AB]$, $[BC]$, $[CD]$, $[DA]$. Usando \omterm{def:g10:coordgeom:system}{coordenadas}
$A(x_A, y_A)$ etc., mostre que $\vect{IJ} = \vect{LK}$ e conclua que os
\omterm{prop:g10:coordgeom:midpoint}{pontos médios} dos lados de \emph{qualquer} quadrilátero formam um
paralelogramo.
\end{exercise}

\section{Problema: A álgebra das setas}

\begin{problem}[{Problema de fim de semana --- ginástica de Chasles, o
baricentro redemonstrado em duas linhas e forças em equilíbrio: para
que servem realmente os vetores}]
\label{pb:g10:vectors:1}
O ano 8 provou que as medianas de um triângulo se encontram a dois
terços do seu comprimento --- com um engenhoso paralelogramo
escondido dentro do triângulo (o problema de fim de semana sobre o
baricentro, no volume do ensino fundamental). Os \omterm{def:g10:vectors:vector}{vetores}
redemonstram isso em duas linhas e ainda dão a concorrência de
brinde. É para isso que serve essa nova álgebra: teoremas viram
contas com setas. Este problema treina a conta (Chasles acima de
tudo), entrega a demonstração de duas linhas, estende Varignon
(\cref{exo:g10:vectors:10}) e termina onde os \omterm{def:g10:vectors:vector}{vetores} nasceram: nas
forças e nas velocidades.

\medskip
\noindent\textbf{Parte I --- Ginástica de Chasles.}
\begin{enumerate}
  \item Simplifique usando a relação de Chasles
        (\cref{thm:g10:vectors:chasles}):
        $\vect{AB} + \vect{BC} + \vect{CD}$; \quad
        $\vect{MA} - \vect{MB}$; \quad
        $\vect{AB} + \vect{CD} + \vect{BC} + \vect{DA}$.
  \item O truque da origem: mostre que, para \emph{qualquer} ponto
        $O$, $\vect{AB} = \vect{OB} - \vect{OA}$.
  \item Demonstre as duas caracterizações do \omterm{prop:g10:coordgeom:midpoint}{ponto médio} $I$ de
        $[AB]$:
        \[
        \vect{IA} + \vect{IB} = \vec 0
        \qquad\text{e}\qquad
        \vect{OI} = \tfrac12\left(\vect{OA} + \vect{OB}\right)
        \ \text{para todo } O .
        \]
  \item Usando $\vect{AB} = \vect{DC}$, reencontre o quarto
        vértice $D$ do paralelogramo $ABCD$ com $A(-1, 2)$,
        $B(3, 4)$, $C(6, 0)$ --- e compare com a resposta obtida
        pelo truque do \omterm{prop:g10:coordgeom:midpoint}{ponto médio} no
        \cref{pb:g10:coordgeom:1}.
  \item Somar um \omterm{def:g10:vectors:vector}{vetor} fixo $\vec u$ a todo ponto do plano
        realiza que transformação --- aquela descoberta com dois
        espelhos no problema de fim de semana sobre os dois
        espelhos e com duas meias voltas no problema de fim de
        semana sobre meias voltas, ambos do volume do ensino
        fundamental? Calcule a \omterm{def:g10:functions:function}{imagem} de $(1, 2)$ por
        $\vec u = (3, -1)$.
\end{enumerate}

\medskip
\noindent\textbf{Parte II --- O baricentro, terceira demonstração.}
Defina o ponto $G$ de um triângulo $ABC$ pela elegante condição
\[
\vect{GA} + \vect{GB} + \vect{GC} = \vec 0 .
\]
\begin{enumerate}[resume]
  \item Usando o truque da origem, mostre que essa condição
        determina $G$ de modo único:
        $\vect{OG} = \frac13\left(\vect{OA} + \vect{OB} +
        \vect{OC}\right)$ para todo $O$ --- de modo que as
        \omterm{def:g10:coordgeom:system}{coordenadas} de $G$ são as \emph{médias} das dos vértices
        (o problema de fim de semana sobre o baricentro, no
        volume do ensino fundamental, viu isso numericamente).
  \item Demonstre que $\vect{AG} = \frac13\left(\vect{AB} +
        \vect{AC}\right)$ e, escrevendo $K$ para o \omterm{prop:g10:coordgeom:midpoint}{ponto médio} de
        $[BC]$, que $\vect{AK} = \frac12\left(\vect{AB} +
        \vect{AC}\right)$. Conclua:
        $\vect{AG} = \frac23\,\vect{AK}$ --- o teorema dos dois
        terços, redemonstrado.
  \item Por que as medianas que partem de $B$ e de $C$ passam
        pelo \emph{mesmo} $G$ sem nenhum cálculo adicional?
        Compare as três demonstrações disponíveis desse teorema
        --- o paralelogramo escondido do ano 8, as \omterm{def:g10:coordgeom:system}{coordenadas} e
        os \omterm{def:g10:vectors:vector}{vetores} --- em uma frase cada.
  \item Verificação numérica: $A(1, 1)$, $B(5, 2)$, $C(3, 6)$.
        Calcule $G$, depois $K$, e verifique
        $\vect{AG} = \frac23 \vect{AK}$ com o critério de
        colinearidade (\cref{thm:g10:vectors:det}).
  \item A física lê $G$ como o ponto de equilíbrio de três massas
        iguais nos vértices. Dobre a massa em $A$ (massas
        $2, 1, 1$): o ponto de equilíbrio passa a ser
        $\vect{OG'} = \frac{2\vect{OA} + \vect{OB} +
        \vect{OC}}{4}$. Calcule $G'$ para o triângulo da questão
        9 e compare sua posição com a de $G$.
\end{enumerate}

\medskip
\noindent\textbf{Parte III --- A colinearidade em ação.}
\begin{enumerate}[resume]
  \item $\vec u(3, -2)$ e $\vec v(-6, 4)$ são \omterm{def:g10:vectors:collinear}{colineares}? E
        $\vec w(2, 5)$, $\vec z(4, 9)$? Decida com o critério do
        determinante.
  \item Os pontos $A(1, 2)$, $B(3, 5)$, $C(7, 11)$ são
        \omterm{def:g10:vectors:collinear}{colineares}?
  \item Para além de Varignon (\cref{exo:g10:vectors:10}): em
        qualquer quadrilátero $ABCD$, as \emph{bimedianas} unem
        os \omterm{prop:g10:coordgeom:midpoint}{pontos médios} de lados opostos ($[IK]$ e $[JL]$).
        Usando \omterm{def:g10:vectors:vector}{vetores} posição ($\vect{OI} = \frac12(\vect{OA} +
        \vect{OB})$ etc.), mostre que as duas bimedianas têm o
        mesmo \omterm{prop:g10:coordgeom:midpoint}{ponto médio}, de \omterm{def:g10:vectors:vector}{vetor} posição
        $\frac14\left(\vect{OA} + \vect{OB} + \vect{OC} +
        \vect{OD}\right)$: o baricentro do quadrilátero.
  \item Tales em uma linha: $M$ e $N$ são os pontos definidos por
        $\vect{AM} = \frac13\vect{AB}$ e
        $\vect{AN} = \frac13\vect{AC}$. Calcule $\vect{MN}$ em
        \omterm{def:g10:functions:function}{função} de $\vect{BC}$ e conclua o paralelismo e a razão.
  \item Generalize a questão 14 para uma razão $k$ qualquer e
        diga em uma frase como o cálculo vetorial absorve de uma
        só vez o teorema da base média e o de Tales.
\end{enumerate}

\medskip
\noindent\textbf{Parte IV --- Forças e velocidades.}
\begin{enumerate}[resume]
  \item Um barco é empurrado por uma correnteza de $3$ km/h para
        leste enquanto rema a $4$ km/h para o norte. Dê o \omterm{def:g10:vectors:vector}{vetor}
        velocidade resultante, seu \omterm{def:g10:numbers:abs}{módulo} e descreva a rota
        efetiva. (Aparece um trio muito antigo.)
  \item Três forças agem sobre um anel: $\vec u(2, 1)$,
        $\vec v(-3, 2)$, $\vec w(1, -3)$. Calcule a resultante. O
        que o anel faz? E por que a \omterm{def:g10:vectors:sum}{soma} das setas ao longo de um
        polígono \emph{fechado} é sempre $\vec 0$ (Chasles)?
  \item Duas forças $\vec u(4, -1)$ e $\vec v(-1, 3)$ agem sobre
        um gancho. Que terceira força mantém o gancho em
        equilíbrio?
  \item Um nadador atravessa um rio apontando perpendicularmente
        à margem a $2$ m/s; a correnteza corre a $1.5$ m/s. Dê a
        velocidade resultante e seu \omterm{def:g10:numbers:abs}{módulo}. O rio tem $50$ m de
        largura: quanto tempo dura a travessia e a que distância
        rio abaixo o nadador chega?
  \item Final: um objeto, três trajes --- seta de deslocamento,
        par de \omterm{def:g10:coordgeom:system}{coordenadas}, força ou velocidade. Liste os três
        teoremas clássicos que este problema redemonstrou com
        álgebra vetorial e nomeie a única grandeza geométrica que
        os \omterm{def:g10:vectors:vector}{vetores} deste capítulo ainda não conseguem medir --- a
        ferramenta que resolve isso chega com o produto escalar,
        no ano 11.
\end{enumerate}
\end{problem}
